\section{Deriving the Derivatives (Backward Pass)}

In this section, the analytical expression for the Jacobian Matrix,
\begin{equation*}
    J(\mathbf{x}):=\frac{\partial F(\mathbf{x},\mathcal{M}_0)}{\partial \mathbf{x}} \in \mathbb{R}^{2m\times 2n},
\end{equation*}
will be derived. The Jacobian Matrix will be later used to solve \cref{eq:main_function}, as discusses in \autoref{sec:solving_eq}.\\

To simplify notation, let
\begin{equation*}
    F_h:= h_n - \theta L_{n+1} q_n^2,\quad F_q := q_n,
\end{equation*}
which are simplified terms of \cref{eq:F_estimate}.

The Jacobian can then be constructed by solving 
$\frac{\partial F_h}{\partial L_i}$, 
$\frac{\partial F_h}{\partial c_i}$,
$\frac{\partial F_q}{\partial L_i}$ and
$\frac{\partial F_q}{\partial c_i}$
for $i\in \{1,2,\dots,n\}$.
Given that
\begin{equation*}
    \frac{\partial}{\partial L_i} L_{n+1} = \frac{\partial}{\partial L_i} \left(L-\sum_{i=1}^n L_i\right) = -1,
\end{equation*}
the derivatives equate to
\begin{align}
    \frac{\partial F_h}{\partial L_i} &= \frac{\partial h_n}{\partial L_i} 
    - 2 \theta L_{n+1} q_n \frac{\partial q_{n}}{\partial L_i} 
    + \theta q_n^2 \label{eq:partial_fh}\\
    \frac{\partial F_h}{\partial c_i} &= \frac{\partial h_n}{\partial c_i}
     - 2 \theta L_{n+1} q_n \frac{\partial q_{n}}{\partial c_i}\\
    \frac{\partial F_q}{\partial L_i} &= \frac{\partial q_{n}}{\partial L_i}\\
    \frac{\partial F_q}{\partial c_i} &= \frac{\partial q_{n}}{\partial c_i} \label{eq:partial_fq}
\end{align}
The remaining differential equations,
\begin{equation*}
    \frac{\partial h_{n}}{\partial L_i},\; 
    \frac{\partial h_{n}}{\partial c_i},\; 
    \frac{\partial q_{n}}{\partial L_i},\; 
    \frac{\partial q_{n}}{\partial c_i},
\end{equation*}
can be expressed in matrix form,
\begin{equation*}
\frac{\partial \phi_{n}}{\partial \mathbf{x}_i} = 
    \begin{bmatrix}
        \frac{\partial h_n}{\partial L_i} & \frac{\partial h_n}{\partial c_i} \\
        \frac{\partial q_n}{\partial L_i} & \frac{\partial q_n}{\partial c_i}
    \end{bmatrix},
\end{equation*}
and solved using the chain rule by
\begin{equation*}
    \frac{\partial \phi_{n}}{\partial \mathbf{x}_i} = 
    \frac{\partial \phi_{n}}{\partial \phi_{i}} \frac{\partial \phi_i}{\partial \mathbf{x}_i} =
    \frac{\partial \phi_{n}}{\partial \phi_{n-1}} \cdots
    \frac{\partial \phi_{i+1}}{\partial \phi_{i}}
    \frac{\partial \phi_{i}}{\partial \mathbf{x}_i}.
\end{equation*}
Let
\begin{equation*}
    A_i:=\frac{\partial \phi_i}{\partial \phi_{i-1}} = \frac{\partial f_{fp}(\phi_{i-1}; \mathbf{x}_i )}{\partial \phi_{i-1}},
\end{equation*}
be the local Jacobian for a single forward pass iteration, where $\phi_{i-1}$ are treated as independent state variables 
\footnote{Treating $\phi_{i-1}$ as independent state variables means that $\frac{\partial q_{i-1}}{\partial h_{i-1}}=0$ even though $q_{i-1}$ is a function of $h_{i-1}$ in the global view.}. From \autoref{lem:forwardpass}, it follows that
\begin{equation*}
    \frac{\partial \phi_k}{\partial \phi_i}=\mathbf{0}, k<i.
\end{equation*}

From \cref{eq:upstream_head}, the partial derivatives for $h_i$ in $A_i$ are computed as
\begin{align*}
    \frac{\partial h_i}{\partial h_{i-1}} &= 1, \\
    \frac{\partial h_i}{\partial q_{i-1}} &= -2 \theta L_i q_{i-1}.
\end{align*}
From \cref{eq:upstream_flow}, the partial derivatives of $q_i$ are
\begin{align*}
    \frac{\partial q_i}{\partial h_{i-1}} &= \frac{\partial q_i}{\partial h_{i}}\frac{\partial h_i}{\partial h_{i-1}} = -\frac{c_i}{2 \sqrt{h_{i}}}, \\
    \frac{\partial q_i}{\partial q_{i-1}} &= 1 - \frac{c_i}{2 \sqrt{h_i}}\frac{\partial h_i}{\partial q_{i-1}}
    = 1+\frac{c_i \theta L_i q_{i-1}}{\sqrt{h_i}}.
\end{align*}
Put in matrix form, $\frac{\partial \phi_i}{\partial \phi_{i-1}}$, is expressed as 
\begin{equation}
\label{eq:A_i}
A_i=\begin{bmatrix}
    1 & -2 \theta L_i q_{i-1} \\[0.4em]
    \displaystyle
    -\frac{c_i}{2 \sqrt{h_{i}}} & \displaystyle 1+\frac{c_i \theta L_i q_{i-1}}{\sqrt{h_i}}
\end{bmatrix}.
\end{equation}

The partial derivatives in $A_i$ can be chained for the backwards pass by defining
\begin{equation*}
    G_i := \frac{\partial \phi_n}{\partial \phi_i} = A_n A_{n-1} \cdots A_{i+1},
\end{equation*}
which can be solved through backward iteration by
\begin{equation*}
    G_i = G_{i+1}A_{i+1}, \quad G_n:=\mathbf{I}.
\end{equation*}

Let
\begin{equation*}
    B_i:=\frac{\partial \phi_i}{\partial \mathbf{x}_i}.
\end{equation*}
From \autoref{lem:forwardpass}, it follows that
\begin{equation*}
    \frac{\partial \phi_k}{\partial \mathbf{x}_i} = \mathbf{0},\; k<i.
\end{equation*}
The differential equations of $B_i$ can then be solved from \cref{eq:upstream_head,eq:upstream_flow},
\begin{align*}
    \frac{\partial h_i}{\partial L_i} &= -\theta q_{i-1}^2,\\
    \frac{\partial h_i}{\partial c_i} &= 0, \\
    \frac{\partial q_i}{\partial L_i} &= - \frac{c_i}{2\sqrt{h_i}} \frac{\partial h_i}{\partial L_i} = 
    \frac{c_i \theta q_{i-1}^2}{2 \sqrt{h_i}}, \\
    \frac{\partial q_i}{\partial c_i} &= -\sqrt{h_i}, \\
\end{align*}
which yields
\begin{equation}
\label{eq:B_i}
B_i=\begin{bmatrix}
    -\theta q_{i-1}^2 & 0 \\[0.4em]
    \displaystyle
    \frac{c_i \theta q_{i-1}^2}{2 \sqrt{h_i}} & -\sqrt{h_i}
\end{bmatrix}.
\end{equation}
To summarise, the partial derivatives are derived by solving
\begin{equation}
    \label{eq:derivative_equation}
    \frac{\partial \phi_{n}}{\partial \mathbf{x}_i} =
     A_n A_{n-1} \cdots A_{i+1} B_i =
    G_i B_i.
\end{equation}

A proposed method for calculating $\phi_i$ and $\frac{\partial \phi_n}{\partial \mathbf{x}_i}$ for all $i\in\{1,2,\dots,n\}$ in $O(n)$ time is given by \autoref{alg:forward_backwards_pass}. The returned list $D$ takes the form

\begin{equation*}
    D = \left\{\frac{\partial \phi_n}{\partial \mathbf{x}_1}, \cdots ,\frac{\partial \phi_n}{\partial \mathbf{x}_n} \right\}
\end{equation*}

\begin{algorithm}[ht]
\caption{Forward and Backwards Pass}
\label{alg:forward_backwards_pass}
%\caption{hello}

\DontPrintSemicolon
\SetKwFunction{func}{ForwardBackwardPass}
\SetKwFunction{Afunc}{CalcA}
\SetKwFunction{Bfunc}{CalcB}
\SetKwFunction{Reverse}{Reverse}
\SetKwFunction{return}{Return}
\SetKwProg{Fn}{Function}{:}{}
\SetKwInput{KwNote}{\normalfont Note}
\Fn{\func{$L,c,h_0,q_0,n,\theta$}}{
    \KwNote{\textit{$L[0]$ and $c[0]$ are dummy values to make indexing consistent with equations: ($L_i:=L[i], c_i:=c[i]$).}}
    \;
    $h \gets [h_0], \; q \gets [q_0]$\;
    \tcp{Forward pass}
    \For{$i \gets 1$ \KwTo n}{
        $h.\mathrm{append}\!\left(h[i-1] - \theta \cdot L[i] \cdot q[i-1]^2\right)$ \tcp*[r]{\cref{eq:upstream_head}}

        $q.\mathrm{append}\!\left(q[i-1] - c[i] \cdot \sqrt{h[i]}\right)$ \tcp*[r]{\cref{eq:upstream_flow}}
    }
    \tcp{Backward pass}
    G $\gets I^{2\times2}$,$\;$ D $\gets$ [ ] \;
    \For{$i \gets n$ \KwTo $1$}{
        B $\gets$ \Bfunc{$c[i],h[i],q[i-1],\theta$} \tcp*[r]{\cref{eq:B_i}}
        $\text{D}.\mathrm{append}\!\left( G\, @\, B\right)$ \;
        \If{$i > 1$}{
            A $\gets$ \Afunc{$L[i],c[i],h[i],q[i-1],\theta$} \tcp*[r]{\cref{eq:A_i}}
            G $\gets$ G $@$ A
        }
    }
    D $\gets$ \Reverse{\normalfont D} \;
    \return $h$, $q$, D \; 
}
\end{algorithm}

The full Jacobian Matrix can be constructed by some simple rearranging. Let 
\begin{align*}
    J^{(j)}(\mathbf{x)} :=& \frac{\partial F(\mathbf{x},\mu_0^{(j)})}{\partial \mathbf{x}} \in \mathbb{R}^{2\times 2n}\\[0.2em]
    =& \begin{bmatrix}
    \frac{\partial F_h^{(j)}}{\partial L_1} & \cdots & \frac{\partial F_h^{(j)}}{\partial L_n} & 
    \frac{\partial F_h^{(j)}}{\partial c_1} & \cdots & \frac{\partial F_h^{(j)}}{\partial c_n} \\[0.4em]
    \frac{\partial F_q^{(j)}}{\partial L_1} & \cdots & \frac{\partial F_q^{(j)}}{\partial L_n} & 
    \frac{\partial F_q^{(j)}}{\partial c_1} & \cdots & \frac{\partial F_q^{(j)}}{\partial c_n}
    \end{bmatrix},
\end{align*}
where all partial derivatives are defined by adding the superscript $(j)$ to Equations (\ref{eq:partial_fh} $-$ \ref{eq:partial_fq}) and solved by the elements in $D^{(j)}$.

Finally,
\begin{equation}
\label{eq:jacobian}
    J(\mathbf{x}) = \begin{bmatrix}
        J^{(1)}\\
        \vdots\\
        J^{(m)}
    \end{bmatrix}
    \in \mathbb{R}^{2m\times2n}.
\end{equation}